% Part: second-order-logic
% Chapter: syntax-and-semantics

\documentclass[../../../include/open-logic-chapter]{subfiles}

\begin{document}

\olchapter{sol}{syn}{विन्यासमीमांसा आणि चिन्हार्थमीमांसा}

\begin{editorial}
आतापर्यंत येथे
द्वितीय-क्रम
तर्कशास्त्राची
मूलभूत विन्यासमीमांसा
आणि चिन्हार्थमीमांसा
घेतली आहे.
स्वतंत्र प्रकरण
म्हणून हे फारच
लहान आहे.
द्वितीय-क्रम
चलांसाठीचे आदेशन
घेतल्याशिवाय
द्वितीय-क्रम
तर्कशास्त्राच्या
!!{derivation} प्रणालींबद्दल
बोलता येणार नाही;
आणि तेथे काही
सूक्ष्म प्रश्नही
उद्भवतात.
\end{editorial}

\olimport{introduction}

\olimport{terms-formulas}

\olimport{satisfaction}

\olimport{semantic-notions}

\olimport{expressive-power}

\olimport{inf-count}

\OLEndChapterHook

\end{document}
